\Needspace{18\baselineskip}
\section{Recursive deep specimens and their boundary and rooted-path laws}
\label{sec:deep-specimens}
The shallow specimen construction now extends recursively through original
box bodies. It supplies actual filled states and proves the conclusions of
Lemma~26(3), Lemma~39(1), and the jointly rooted context-path consequence for
this deeper representation. These are structural results toward the paper's
specimen argument in \S8.1 and Appendix~F~\cite{paper}. Finite original-hole tuple compatibility and the full concrete safety
definitions are established later in Section~\ref{sec:tuple-safety}. The
complete specimen-interface interpretation remains a separate obligation.

\Needspace{12\baselineskip}
\subsection{Copies with an updated box label}
\begin{definition}[Specified-label copy]\label{term:label-copy}
A specified-label copy is a graph with a distinguished sole edge, a
specified label on that edge, and source and target lists enumerating the
full ordered input and output boundaries. Unlike an unchanged corolla
copy, its label need not be the original host label: it may contain a
new box body. The graph itself still satisfies all finite typed graph
invariants. The literal corolla is such a copy, with its sole edge $0$;
every element of $\mathrm{Fin}(1)$ equals $0$.
\end{definition}
\lean{Graph/DeepFillings}{11}{16}{Linear.LabelCopy}
\lean{Graph/DeepFillings}{18}{20}{Linear.LabelCopy.corolla}

If the specified label is globally focus-free, so is the copy.
Every edge equals its distinguished edge, whose label is the specified
label, reducing the claim to the supplied label predicate.
\lean{Graph/DeepFillings}{22}{27}{Linear.LabelCopy.focusFree}

\begin{lemma}[A non-operation copy is one-way]
If a copy's specified label is not an operation, its graph is one-way.
\end{lemma}
\begin{proof}
Use the characterization by operation reachability of term outputs.
Every output belongs to the sole edge's target list, because that list
enumerates the whole output boundary. If such an output were
operation-reachable, incoming-edge uniqueness would force that edge
to be an operation. Its label equality contradicts the hypothesis.
\end{proof}
\lean{Graph/DeepFillings}{30}{40}{Linear.LabelCopy.oneWay}

\clearpage
\subsection{Recursive paired filling certificates}
\begin{definition}[Layered pre-template]\label{term:layered-template}
A layered pre-template chooses a concrete, interface-indexed pre-template
$T_d$ for every depth bound $d$. Each related pair is globally focus-free.
This definition does not identify graphs or relations at different
depth bounds. Section~\ref{sec:paper-encoding} constructs coherent families
from a single global pre-template; arbitrary families need not be coherent.
\end{definition}
\lean{Graph/DeepFillings}{44}{46}{Linear.LayeredPreTemplate}

\begin{definition}[Recursive paired filling]\label{term:deep-filling}
The predicate $\mathsf{DeepFilling}_T(C,L,R)$ certifies two actual filled
graphs $L,R$ of the same interface and depth as $C$. A node chooses two
outer filling families $H,K$, certifies each original edge as below,
and sets $L=C[H]$ and $R=C[K]$ using the concrete expansion construction.
It is mutually defined with the edge certificates, so every recursive
body filling is itself such a finite construction.
\end{definition}
\lean{Graph/DeepFillings}{49}{57}{Linear.DeepFilling}

\begin{definition}[Recursive edge certificate]\label{term:deep-filling-edge}
An original edge has one of three certificates.
It can be retained by unchanged corolla copies on both sides.
It can have a hole label and be filled by a pair related by $T_d$.
Or it can be a box with bound/free metadata $(b,f)$ and original body
$B$: recursively fill $B$ to bodies $L,R$, and use specified-label
copies of the boxes $(b,f,L)$ and $(b,f,R)$ as the two outer fillings.
The body interfaces remain one term input and $b+f$ term outputs.
Recursion follows original boxes, not the interiors of newly supplied
hole replacements. Retained holes are permitted by this graph relation;
the specimen's final hole-freedom requirement excludes unresolved holes.
\end{definition}
\lean{Graph/DeepFillings}{61}{82}{Linear.DeepFillingEdge}

Every shallow specimen-edge certificate is a recursive edge certificate:
use the retained or filled constructor without descending into any box.
Thus this representation extends the earlier one directly.
\lean{Graph/DeepFillings}{85}{90}{Linear.SpecimenEdge.deep}

\begin{theorem}[Global focus freedom under recursive filling]
If $\mathsf{DeepFilling}_T(C,L,R)$ and $C$ is globally focus-free, then
both $L$ and $R$ are globally focus-free.
\end{theorem}
\begin{proof}
Induct on the depth bound. Each result edge is an edge in one outer
filling, so it suffices to prove focus freedom of every filling.
At depth zero an edge certificate is retained or filled. A retained
copy inherits the original edge's focus freedom; a filled pair obtains
it from the pre-template domain.

At successor depth the same cases apply, together with a box case.
The original box's label is focus-free, so its original body is
globally focus-free. Apply the induction hypothesis to the certified
body filling. The two new box labels are therefore focus-free, and the
specified-label copy lemma gives focus freedom of their outer copies.
This accounts for all descendant edges, not just outer labels.
\end{proof}
\lean{Graph/DeepFillings}{93}{121}{Linear.DeepFilling.focusFree}

For one edge, focus freedom of its original label already suffices
for focus freedom of both fillings. Retained and hole cases are as above;
the box case applies the global recursive theorem to its body and then
the copy lemma.
\lean{Graph/DeepFillings}{124}{134}{Linear.DeepFillingEdge.focusFree}

A shallow focus must take the retained branch of a recursive edge
certificate. Its label cannot equal a hole label, and cannot equal a
box label. Excluding those two cases leaves the pair of unchanged
focus-copy certificates.
\lean{Graph/DeepFillings}{137}{146}{Linear.DeepFillingEdge.focus_copies}

\begin{lemma}[Outer paired path hypothesis]
If $T_d$ is output-closed, the outer families of a recursive filling
satisfy the paired expansion path hypothesis.
\end{lemma}
\begin{proof}
A retained edge supplies its two unchanged copies.
At a hole, output closure gives a one-way filling on at least one side,
with the choice allowed to vary by edge. At a rebuilt box, its specified
label is still a box and hence is not an operation. The left copy is
therefore one-way by the non-operation copy lemma (as is the right copy,
although one side suffices). This argument needs output closure only
for the outer layer's related hole fillings. Inner focus freedom was
proved recursively; inner output closure is unnecessary for this
outer path obstruction.
\end{proof}
\lean{Graph/DeepFillings}{150}{158}{Linear.DeepFillingEdge.outputClosed}

\Needspace{12\baselineskip}
\subsection{Actual deep specimen states}
\begin{definition}[Concrete recursive specimen]\label{term:deep-specimen}
A recursive specimen consists of a unique-hole context $C$ with outer
interface $[\mathsf{term}]\to[]$, a distinguished shallow focus, an
allowed erased context, two outer filling families, and a recursive
edge certificate at every original outer edge. Both resulting
expansions must be globally hole-free. These last fields certify that
the constructed results are states; the record contains no assumed
boundary agreement, rootedness obstruction, or path-reflection law.
\end{definition}
\lean{Graph/DeepSpecimens}{11}{24}{Linear.DeepSpecimen}

Its edge certificates assemble into a recursive paired filling.
At the focus edge, the preceding retained-focus theorem supplies a
pair of copy certificates; classical choice selects one such pair
for the subsequent state constructors.
\lean{Graph/DeepSpecimens}{30}{31}{Linear.DeepSpecimen.filling}
\lean{Graph/DeepSpecimens}{33}{35}{Linear.DeepSpecimen.focus_copies}
\lean{Graph/DeepSpecimens}{37}{37}{Linear.DeepSpecimen.focusCopies}

Every other original edge is globally focus-free by uniqueness of
the original focus. Its recursive edge certificate therefore gives
globally focus-free fillings on both sides. The expansion focus
construction can now retain the original focus with its mode and
mapped source and target. Together with the specimen's two closedness
certificates this constructs actual left and right states.
\lean{Graph/DeepSpecimens}{39}{41}{Linear.DeepSpecimen.other_free}
\lean{Graph/DeepSpecimens}{43}{44}{Linear.DeepSpecimen.leftState}
\lean{Graph/DeepSpecimens}{46}{47}{Linear.DeepSpecimen.rightState}

Output closure supplies the outer paired path hypothesis edge by edge,
using the recursive edge lemma just proved.
\lean{Graph/DeepSpecimens}{51}{53}{Linear.DeepSpecimen.outputClosedPair}

\begin{theorem}[Deep specimen obstruction: Lemma 26(3)]
Suppose the outer pre-template is output-closed. If an original context
path ends at the focus source and contains a non-operation edge, at
least one of the two constructed states is not rooted.
\end{theorem}
\begin{proof}
Apply the checked paired expansion obstruction to the two families.
The retained focus copies and all other focus-free fillings have just
been derived, the states are closed by their certificates, and the
paired path hypothesis follows from output closure and the treatment
of rebuilt boxes. If both states were rooted, their focus sources
would be operation-reachable. Paired path reflection would make every
edge of the original path an operation, contradicting the chosen edge.
The original path is an outer context path; the fillings may occur
arbitrarily deeply inside its box labels.
\end{proof}
\lean{Graph/DeepSpecimens}{57}{65}{Linear.DeepSpecimen.not_both_rooted}

Equivalently, under joint rootedness every edge of any context path
to the focus source is an operation. The direct proof applies paired
path reflection to the two operation-reachability witnesses obtained
from rootedness, using that the focus source has term type. This is
the rooted-context path consequence needed in the later locality
argument, not the entire entering/exiting classification of Lemma~39.
\lean{Graph/DeepSpecimens}{68}{76}{Linear.DeepSpecimen.rooted_context_path}

\begin{theorem}[Deep specimen boundary agreement: Lemma 39(1)]
The two actual states are initial simultaneously and final simultaneously.
\end{theorem}
\begin{proof}
Expansion retains the focus mode and preserves whether its source
is an input. Thus each expanded focus is initial exactly when the
original focus is initial, and final exactly when the original focus
is final. Compose the left equivalences with the inverses of the
right equivalences. Output closure and rootedness are not required.
\end{proof}
\lean{Graph/DeepSpecimens}{79}{83}{Linear.DeepSpecimen.boundary_agreement}

Each result graph also admits a finite ordinary shallow-substitution
realization of its outer filling family, up to graph isomorphism.
Apply finite sequential realization to the left family and then the
right. These runs may rebuild a whole outer box edge from its already
filled body. Consequently this result alone does not identify the recursive
representation with the paper's original-hole-only tuple substitution. That
separate adequacy theorem is proved in Section~\ref{sec:tuple-safety}.
\lean{Graph/DeepSpecimens}{87}{94}{Linear.DeepSpecimen.outer_realizations}

To embed an earlier shallow specimen, retain its context, focus,
allowedness certificate, two families, and closedness proofs, and
convert each shallow edge certificate to a recursive one.
\lean{Graph/DeepSpecimens}{99}{109}{Linear.ShallowSpecimen.toDeep}

\Needspace{12\baselineskip}
\subsection{A genuinely deep closed-state example}
Use two Boolean-named operations. Both have no term arguments; the
true operation exposes one thunk of binder arity zero, while the false
operation exposes no outputs. Let the unique hole have name $7$ and
one term input and output. Its corolla is the original box body.
The replacement is a unary variable corolla.
\lean{Graph/DeepSpecimenExamples}{8}{8}{Linear.DeepSpecimenExamples.ops}
\lean{Graph/DeepSpecimenExamples}{9}{9}{Linear.DeepSpecimenExamples.hole}
\lean{Graph/DeepSpecimenExamples}{10}{11}{Linear.DeepSpecimenExamples.body}
\lean{Graph/DeepSpecimenExamples}{12}{13}{Linear.DeepSpecimenExamples.wire}

At every depth, the example pre-template relates equal graphs that
are focus-free and one-way. Equality transfers focus freedom to the
second member; output closure follows from the supplied one-wayness
of the first. The variable wire is focus-free and hole-free by its
label, and is one-way by the non-operation copy lemma.
\lean{Graph/DeepSpecimenExamples}{15}{17}{Linear.DeepSpecimenExamples.template}
\lean{Graph/DeepSpecimenExamples}{19}{19}{Linear.DeepSpecimenExamples.template_closed}
\lean{Graph/DeepSpecimenExamples}{21}{21}{Linear.DeepSpecimenExamples.wire_free}
\lean{Graph/DeepSpecimenExamples}{22}{22}{Linear.DeepSpecimenExamples.wire_closed}
\lean{Graph/DeepSpecimenExamples}{23}{24}{Linear.DeepSpecimenExamples.wire_oneWay}

The sole body edge is assigned the variable wire. Expand this family
to obtain the filled body. The hole-label equation and related pair
give a recursive filling certificate. The resulting body is hole-free
because all its inserted graphs are the hole-free variable wire.
\lean{Graph/DeepSpecimenExamples}{26}{26}{Linear.DeepSpecimenExamples.bodyFill}
\lean{Graph/DeepSpecimenExamples}{27}{27}{Linear.DeepSpecimenExamples.filledBody}
\lean{Graph/DeepSpecimenExamples}{29}{30}{Linear.DeepSpecimenExamples.body_filling}
\lean{Graph/DeepSpecimenExamples}{32}{33}{Linear.DeepSpecimenExamples.filledBody_closed}

The outer graph has three edge names: maker, box, and stop. Its two
internal vertices have thunk and term type, respectively. The maker
runs from the unique term input to the thunk vertex; the box runs from
that vertex to the term vertex; the stop operation consumes the term
vertex and has no outputs. The box has no bound variables, one free
term port, and the original hole corolla as its body.
Finite enumerations list the two vertices and three edges.
All nonempty incidence lists are singletons. Case analysis on the
three edges proves source and target uniqueness. Each internal vertex
has exactly the required incoming and outgoing occurrence; the input
belongs to the maker's source, and no outputs exist. These observations
prove coverage. The typing equations hold by the definitions of the
two operations, the box interface, and the vertex labels.
\lean{Graph/DeepSpecimenExamples}{36}{37}{Linear.DeepSpecimenExamples.Part}
\lean{Graph/DeepSpecimenExamples}{39}{83}{Linear.DeepSpecimenExamples.outer}

The outer graph is globally focus-free: its operation labels contain
no focus, and the box's body is a hole corolla. Insert a search focus
at the term vertex immediately after the box, using the existing
verified focus-insertion construction. The vertex has term type and
is not an output, so the construction gives a well-formed focussed
context.
\lean{Graph/DeepSpecimenExamples}{85}{90}{Linear.DeepSpecimenExamples.outer_free}
\lean{Graph/DeepSpecimenExamples}{92}{93}{Linear.DeepSpecimenExamples.context}
\lean{Graph/DeepSpecimenExamples}{95}{96}{Linear.DeepSpecimenExamples.focus}

Its only hole occurrence is the sole edge inside the original box.
To prove this, split the outer edge into the three original edges or
the newly inserted focus. Operation and focus labels have neither
shallow holes nor children. At the box, the shallow occurrence is
not a hole; its only inner edge has index $0$. Thus any two hole
occurrences coincide, proving the unique-hole context property.
There is an explicit witness for that deep occurrence, while no
outer edge has a hole label.
\lean{Graph/DeepSpecimenExamples}{99}{120}{Linear.DeepSpecimenExamples.context_hole_only}
\lean{Graph/DeepSpecimenExamples}{122}{125}{Linear.DeepSpecimenExamples.context_unique}
\lean{Graph/DeepSpecimenExamples}{128}{129}{Linear.DeepSpecimenExamples.context_has_deep_hole}
\lean{Graph/DeepSpecimenExamples}{131}{135}{Linear.DeepSpecimenExamples.context_no_shallow_hole}

\begin{lemma}[This context cannot be filled shallowly to a closed graph]
No pair of families certified by the old shallow specimen-edge rules
on this same context has a hole-free left expansion.
\end{lemma}
\begin{proof}
Every shallow filled-hole branch is impossible, since there is no
shallow hole label. Thus every left filling is an unchanged corolla
copy. The identity-copy expansion theorem gives an isomorphism from
that expansion to the original context. Recursive isomorphism
preserves the explicit deep hole occurrence, contradicting claimed
hole freedom of the expansion.
\end{proof}
\lean{Graph/DeepSpecimenExamples}{139}{152}{Linear.DeepSpecimenExamples.no_shallow_closed_filling}

For the new recursive family, retain both operation edges and the
focus by their literal corollas, but replace the outer box by a box
whose body is the already filled body. The recursive box certificate
uses the body-filling proof and the literal specified-label copy
certificates. Retained edges use their corolla identities. Each filling
is hole-free: the rebuilt box inherits its body's hole freedom, and
the operation and focus corollas contain no shallow or deep holes.
The expansion is therefore globally hole-free.
\lean{Graph/DeepSpecimenExamples}{154}{161}{Linear.DeepSpecimenExamples.filling}
\lean{Graph/DeepSpecimenExamples}{163}{177}{Linear.DeepSpecimenExamples.filling_edges}
\lean{Graph/DeepSpecimenExamples}{179}{196}{Linear.DeepSpecimenExamples.filling_closed}

Package these data as a recursive specimen with unrestricted allowed
contexts. Both sides use the displayed family. The original box edge
alone forms a context path to the focus source and is not an operation.
The deep obstruction theorem yields the disjunction of non-rootedness
for every classification of operations. The general boundary theorem
also gives initial/final agreement for the actual constructed states.
\lean{Graph/DeepSpecimenExamples}{198}{206}{Linear.DeepSpecimenExamples.specimen}
\lean{Graph/DeepSpecimenExamples}{209}{215}{Linear.DeepSpecimenExamples.specimen_not_rooted}
\lean{Graph/DeepSpecimenExamples}{217}{220}{Linear.DeepSpecimenExamples.specimen_boundary}

This example distinguishes the new representation from the old one by
a proved impossibility result, not merely by its choice of notation.
Finite tuple compatibility and concrete quasi-specimen restrictions are proved
in Section~\ref{sec:tuple-safety}. The remaining specimen work includes the
locality, safety, and robustness bridges listed in
Appendix~\ref{app:remaining-work}.
